% ============================================================
%  GLM PORTFOLIO -- SHARED PREAMBLE
%  Playful-but-structured design, ECRA phase system
% ============================================================

\usepackage[T1]{fontenc}
\usepackage[utf8]{inputenc}
\usepackage{charter}
\usepackage[scaled=0.92]{helvet}
\renewcommand{\familydefault}{\rmdefault}

\usepackage{amsmath,amssymb,amsthm}
\usepackage{xcolor}
\usepackage[most]{tcolorbox}
\usepackage{geometry}
\usepackage{titlesec}
\usepackage{fancyhdr}
\usepackage{enumitem}
\usepackage{graphicx}
\usepackage{tikz}
\usepackage{booktabs}
\usepackage{listings}
\usepackage[colorlinks=true,linkcolor=tealmain_,citecolor=tealmain_,urlcolor=amberacc_]{hyperref}
\usepackage[round]{natbib}
\usepackage{caption}

\geometry{a4paper, margin=2.6cm, headsep=1.1cm}

% avoid blank pages from \frontmatter/\mainmatter/\part forcing odd-page starts
\let\cleardoublepage\clearpage
\let\cleardoubleoddpage\clearpage
\let\cleardoubleevenpage\clearpage

% ---------- Colour palette ----------
% Core: teal (Cognition), amber (Reflection)
% Extended: slate-blue (Exploration), graphite (Application) -- tints/neutrals, not new loud hues
\definecolor{inkblack}{HTML}{1F2328}
\definecolor{tealmain}{HTML}{146C6E}
\definecolor{tealsoft}{HTML}{E3F2F1}
\definecolor{amberacc}{HTML}{D97C2B}
\definecolor{ambersoft}{HTML}{FBEBDA}
\definecolor{slateacc}{HTML}{3B5A76}
\definecolor{slatesoft}{HTML}{E9EFF4}
\definecolor{graphite}{HTML}{4B4F52}
\definecolor{graphitesoft}{HTML}{EDEDED}
\definecolor{greyline}{HTML}{9AA5A9}
\definecolor{mathsoft}{HTML}{F3F8F8}

% aliases so hyperref (loaded before some defs) can see the colours
\colorlet{tealmain_}{tealmain}
\colorlet{amberacc_}{amberacc}

\color{inkblack}

% ---------- Chapter / section styling ----------
\titleformat{\chapter}[display]
  {\normalfont\sffamily\bfseries}
  {\begin{tikzpicture}[remember picture,overlay]
     \fill[tealmain] (-0.45,-0.15) rectangle (-0.17,0.85);
   \end{tikzpicture}\color{greyline}\Large\MakeUppercase{\chaptertitlename} \thechapter}
  {6pt}
  {\color{tealmain}\Huge}
\titlespacing*{\chapter}{0pt}{10pt}{22pt}

\RedeclareSectionCommand[
  beforeskip=-10pt, afterskip=20pt,
  font=\normalfont\sffamily\bfseries\huge\color{tealmain}
]{part}
\renewcommand{\partformat}{\color{greyline}\sffamily\Large\MakeUppercase{\partname}~\thepart.\quad}

\titleformat{\section}
  {\normalfont\sffamily\bfseries\color{tealmain}\Large}{\thesection}{8pt}{}
\titleformat{\subsection}
  {\normalfont\sffamily\bfseries\color{amberacc}\large}{\thesubsection}{8pt}{}
\titleformat{\subsubsection}
  {\normalfont\sffamily\bfseries\color{amberacc}}{\thesubsubsection}{6pt}{}

\pagestyle{fancy}
\fancyhf{}
\renewcommand{\headrulewidth}{0.4pt}
\fancyhead[LE,RO]{\sffamily\small\thepage}
\fancyhead[RE]{\sffamily\small\nouppercase{\leftmark}}
\fancyhead[LO]{\sffamily\small\nouppercase{\rightmark}}

% ============================================================
%  ECRA PHASE SYSTEM
%  Each phase = a coloured banner heading + a matching box style
% ============================================================

% ---- generic phase-banner heading ----
\newcommand{\phaseheading}[2]{%
  \vspace{14pt}
  \noindent\colorbox{#1}{\parbox{\dimexpr\linewidth-12pt}{\vspace{3pt}\sffamily\bfseries\color{white}\large #2\vspace{3pt}}}%
  \vspace{8pt}\par
}

% ---- EXPLORATION: slate-blue, "browser frame" screenshot box ----
\newcommand{\Exploration}{\phaseheading{slateacc}{Exploration}}

\newenvironment{screenshot}[1][]{%
  \begin{center}
  \begin{tcolorbox}[
    colback=white, colframe=slateacc, boxrule=0.6pt, arc=1pt,
    width=0.92\linewidth,
    top=2pt,bottom=2pt,left=2pt,right=2pt,
    before skip=8pt, after skip=4pt
  ]
  \ifx\\#1\\\else{\sffamily\footnotesize\color{slateacc} #1}\\[2pt]\fi
}{%
  \end{tcolorbox}
  \end{center}
}

% informal exploration-phase source note (not a formal \cite)
\newcommand{\explorationsource}[1]{{\sffamily\footnotesize\color{slateacc} Source: #1}}

% ---- COGNITION: teal, definitions + theorems + formal citation ----
\newcommand{\Cognition}{\phaseheading{tealmain}{Cognition}}

\newtcolorbox{defbox}[1][]{
  colback=tealsoft, colframe=tealmain, boxrule=0.6pt, arc=2pt,
  fonttitle=\sffamily\bfseries, title={Definition#1},
  left=8pt,right=8pt,top=5pt,bottom=5pt
}
\newtheorem{theorem}{Theorem}[chapter]
\newtheorem{lemma}[theorem]{Lemma}
\newtheorem{proposition}[theorem]{Proposition}
\renewcommand{\qedsymbol}{$\blacksquare$}

% glossary-style term heading (for Chapter 3's Language of GLMs)
\newcommand{\term}[1]{\vspace{4pt}\subsubsection*{\sffamily\bfseries\color{tealmain} #1}}

% unnumbered chapter heading (Part IV evaluation topics), same visual family, no "Chapter N" tab
\newcommand{\uchapter}[1]{%
  \clearpage
  \phantomsection
  \vspace*{10pt}
  {\sffamily\bfseries\color{tealmain}\huge #1\par}
  \vspace{18pt}
  \addcontentsline{toc}{chapter}{#1}
  \markboth{#1}{#1}
}

% ---- REFLECTION: amber, question/exercise boxes ----
\newcommand{\Reflection}{\phaseheading{amberacc}{Reflection}}

\newtcolorbox{questionbox}[1][]{
  colback=ambersoft, colframe=amberacc, boxrule=0.6pt, arc=2pt,
  fonttitle=\sffamily\bfseries, title={Question#1},
  left=8pt,right=8pt,top=5pt,bottom=5pt
}

% ---- APPLICATION: graphite/neutral, code ----
\newcommand{\Application}{\phaseheading{graphite}{Application}}

\lstdefinestyle{rcode}{
  backgroundcolor=\color{graphitesoft},
  basicstyle=\ttfamily\footnotesize\color{inkblack},
  keywordstyle=\color{tealmain}\bfseries,
  commentstyle=\color{greyline}\itshape,
  numberstyle=\tiny\color{greyline},
  stringstyle=\color{amberacc},
  numbers=left,
  numbersep=8pt,
  frame=single,
  rulecolor=\color{greyline},
  framerule=0.4pt,
  breaklines=true,
  showstringspaces=false,
  xleftmargin=10pt,
  framexleftmargin=6pt
}
\lstset{style=rcode}

% "to be completed" placeholder, styled quietly so it doesn't compete visually
\newcommand{\tbc}{{\sffamily\itshape\color{greyline} To be completed.}\par\vspace{6pt}}

% ============================================================
%  MATH HIGHLIGHT DESIGN
%  Every standalone display equation gets its own subtle
%  highlighted band, distinguishing math from prose at a glance.
% ============================================================
\usepackage[framemethod=TikZ]{mdframed}
\mdfdefinestyle{mathstyle}{
  backgroundcolor=mathsoft,
  linecolor=tealmain,
  linewidth=0pt,
  leftline=true,
  topline=false,
  bottomline=false,
  rightline=false,
  linewidth=2.5pt,
  innerleftmargin=12pt,
  innerrightmargin=8pt,
  innertopmargin=6pt,
  innerbottommargin=6pt,
  skipabove=8pt,
  skipbelow=8pt
}
\newmdenv[style=mathstyle]{mathbox}
